\section*{Chapitre \ref{ch:rc:multi-curve-and-collateral-discounting} --- Cadre multi-courbes et actualisation par collatéral}

\begin{solution}{exo:rc:multi-curve-and-collateral-discounting:1}
$2.70-2.55 = 15$ points de base
(la \cref{prop:rc:multi-curve-and-collateral-discounting:basisrates}), ajouté à la
jambe overnight~: la partie qui reçoit l'Euribor six mois verse le €STR capitalisé plus
15 points de base.
\end{solution}

\begin{solution}{exo:rc:multi-curve-and-collateral-discounting:2}
L'argument à courbe unique suppose qu'un prêt de six mois coûte autant que deux
prêts de trois mois ou 180 prêts au jour le jour~; un écart persistant de cinquante
points de base entre prêts à terme non garantis et garantis faisait de chaque ténor un
prix différent. Une jambe variable indexée sur un indice ne valait plus le pair sur une courbe
construite à partir d'une autre, et l'actualisation au taux à terme d'une banque n'était plus le
coût d'une position collatéralisée.
\end{solution}

\begin{solution}{exo:rc:multi-curve-and-collateral-discounting:3}
La courbe €STR actualise~; la \omterm{def:rc:multi-curve-and-collateral-discounting:curves}{courbe de projection} Euribor six mois projette. Avec
€STR plus 10 points de base, le \omterm{def:rc:multi-curve-and-collateral-discounting:csa}{taux de rémunération du collatéral} est plus élevé, si bien que la \omterm{def:rc:multi-curve-and-collateral-discounting:curves}{courbe
d'actualisation} est la courbe €STR décalée de 10 points de base (facteurs d'actualisation multipliés par
$e^{-0.001t}$)~; la \omterm{def:rc:multi-curve-and-collateral-discounting:curves}{courbe de projection}, recalibrée sur elle, se déplace légèrement, et
une transaction hors marché voit sa valeur changer.
\end{solution}

\begin{solution}{exo:rc:multi-curve-and-collateral-discounting:4}
L'écart effectif vaut $\max(0,-b(t)) = (15-2t)$ points de base pour $t<7.5$
et nul après~: $\int_0^{7.5}(15-2t)\,dt = 112.5-56.25 = 56.25$
points de base·années. Le rapport des facteurs d'actualisation vaut $e^{-0.005625}=0.99439$, et
$68\,598\,292\times(1-0.99439) = \text{USD}~384\,782$.
\end{solution}

\begin{solution}{exo:rc:multi-curve-and-collateral-discounting:5}
La courbe unique actualise aux taux Euribor, environ 15 points de base au-dessus des
taux €STR auxquels le swap collatéralisé est réellement financé, elle sous-évalue donc
les flux nets futurs du receveur, de EUR~60\,300. Les deux cadres réévaluent par construction le
swap à dix ans au pair, ils concordent donc au taux pair et ne diffèrent que pour les
transactions hors marché, d'un montant qui croît avec
la distance au pair.
\end{solution}

\begin{solution}{exo:rc:multi-curve-and-collateral-discounting:6}
Un payeur de 1.50~\% alors que le taux pair est bien plus bas détient un passif~: sa valeur
vaut $-\text{EUR}~16.70$ millions, et des taux d'actualisation plus bas alourdissent le passif
plus lourd, soit une perte de EUR~73\,738, symétrique de la barre à dix ans. La chambre
de compensation le dédommage~: le client reçoit EUR~73\,738, versé à partir des
gains des receveurs.
\end{solution}

\begin{solution}{exo:rc:multi-curve-and-collateral-discounting:7}
11.76 points de base, contre une base pair à dix ans de 15. La base pair est une moyenne
pondérée par l'annuité des bases forward sur dix ans~; la base forward part de 21.6 et
décroît, si bien que la moyenne à dix ans dépasse la base forward à dix ans. Couvrir une
position sur base à départ différé avec des swaps de base pair exige toute
la structure par terme.
\end{solution}

\begin{solution}{exo:rc:multi-curve-and-collateral-discounting:8}
Le déposant peut changer de collatéral à tout moment sur la vie de la transaction, si bien
que l'option porte sur tout le chemin de la base, et non sur sa valeur du jour~: même si le
dollar est le moins cher aujourd'hui, les dates futures où l'euro est meilleur abaissent la
valeur d'une créance (le taux effectif est le maximum à chaque date), et la volatilité de la
base ajoute une valeur temps. Valoriser sur la courbe en dollars ignore ces deux effets.
\end{solution}

\begin{solution}{pb:rc:multi-curve-and-collateral-discounting:1}
\textbf{1.} Un actif~: il reçoit 1.50~\% alors que le taux pair est de 0.08~\%, si bien que ses
flux nets futurs sont importants.
\textbf{2.} EUR~145\,577\,447.
\textbf{3.} Les versements fixes et la période Euribor en cours sont connus~; les fixings
Euribor ultérieurs sont projetés à partir de la \omterm{def:rc:multi-curve-and-collateral-discounting:curves}{courbe de projection}.
\textbf{4.} Il a reçu une marge de variation égale à la valeur du swap et verse dessus
du \omterm{def:rc:multi-curve-and-collateral-discounting:pai}{price alignment interest} (PAI) au taux EONIA avant le changement, puis au €STR (8.5
points de base de moins) après~; avec des taux négatifs, le «~paiement~» est une réception.
\textbf{5.} La position du fonds est financée par la marge de variation au taux PAI, si bien que
sa valeur s'accroît à ce taux
(la \cref{prop:rc:multi-curve-and-collateral-discounting:csa}).
\textbf{6.} Chaque facteur d'actualisation est multiplié par $e^{0.00085\,t}$~: 1.0171 à
vingt ans.
\textbf{7.} EUR~146\,774\,873.
\textbf{8.} $+\text{EUR}~1\,197\,426$~: le fonds gagne, car un taux d'actualisation plus bas
augmente la valeur des flux nets.
\textbf{9.} Le fonds verse EUR~1\,197\,426 en espèces à la chambre de compensation, qui
le reverse aux membres dont les positions ont perdu.
\textbf{10.} Les cotations de marché des swaps Euribor n'ont pas bougé pendant le week-end~;
la méthode isole le pur effet d'actualisation, est simple à reproduire par chaque membre,
et rejoint le point de vue du groupe de travail selon lequel les forwards ne dépendent pas
du CSA.
\textbf{11.} EUR~141\,575 par point de base.
\textbf{12.} C'était le risque d'actualisation EONIA, c'est désormais le risque d'actualisation €STR.
EONIA valait €STR plus 8.5 points de base fixes, si bien que les deux risques étaient identiques
et rien n'était à couvrir~; les federal funds et le SOFR évoluent de façon indépendante, si bien que
le passage au dollar a modifié le risque des membres, que des swaps de base ont rétabli.
\textbf{13.} Perdant~: son passif a augmenté à mesure que les facteurs d'actualisation montaient.
\textbf{14.} Le gain vaut à peu près la valeur multipliée par la duration et par 8.5 points de base,
et la valeur d'un swap dans la monnaie comme sa duration croissent avec
l'échéance~; le gain à trente ans vaut 4.1 fois celui à quinze ans.
\textbf{15.} Non~: la compensation le ramène à sa situation initiale. Ses comptes peuvent
afficher un gain et un versement en espèces égal, ainsi qu'un risque d'actualisation légèrement différent.
\textbf{16.} Une actualisation mixte au sein d'un même pool de compensation romprait
l'équivalence entre les positions des membres et la marge de la chambre de compensation, et une
date unique permet à la compensation de compenser sur l'ensemble des comptes.
\textbf{17.} Un amendement bilatéral du CSA au €STR (plus ou moins un spread),
négocié transaction par transaction, avec son propre versement de compensation.
\textbf{18.} Peu de chose~: EONIA valait mécaniquement €STR plus 8.5 points de base, donc le
passage était un transfert déterministe connu, compensé en espèces~; il n'y avait aucune
inconnue à échanger.
\textbf{19.} Résultat nommé~: \emph{week-end du passage} a relevé la valeur du swap du fonds de
EUR~1\,197\,426, qu'il a restitué en compensation en espèces~; après le passage, sa sensibilité à l'actualisation est de
EUR~141\,575 par point de base d'€STR.
\textbf{20.} Le taux versé sur le collatéral qui le garantit.
\end{solution}

\begin{solution}{iq:rc:multi-curve-and-collateral-discounting:1}
Parce que la transaction est financée par son collatéral~: le détenteur d'une transaction de valeur positive
détient un collatéral en espèces égal à sa valeur et verse sur ce montant le taux au jour le jour, de sorte que
la valeur croît à ce taux. L'indice ne détermine que les flux de trésorerie,
qui sont projetés sur sa propre courbe.
\iqlookfor{actualisation liée au taux du collatéral, projection sur l'indice.}
\end{solution}

\begin{solution}{iq:rc:multi-curve-and-collateral-discounting:2}
D'abord la \omterm{def:rc:multi-curve-and-collateral-discounting:curves}{courbe d'actualisation} au taux au jour le jour (€STR), issue des OIS~; ensuite, celle-ci étant fixée,
la \omterm{def:rc:multi-curve-and-collateral-discounting:curves}{courbe de projection} de l'Euribor à six mois, construite à partir du fixing à six mois, des
contrats FRA ou contrats à terme, et des swaps contre l'Euribor à six mois, chacun
évalué avec les taux forward de l'Euribor actualisés sur l'€STR. La \omterm{def:rc:multi-curve-and-collateral-discounting:curves}{courbe d'actualisation} vient en premier
parce que le prix de chaque instrument de projection en dépend, et non l'inverse.
\iqlookfor{l'ordre et la raison.}
\end{solution}

\begin{solution}{iq:rc:multi-curve-and-collateral-discounting:3}
Deux jambes variables de ténors différents (Euribor à six mois contre €STR ou
Euribor à trois mois), avec un spread sur l'une pour la rendre équitable. Il est positif
parce qu'un prêt non garanti à six mois porte une prime de crédit et de liquidité plus élevée
qu'une succession roulante de prêts plus courts ou garantis.
\iqlookfor{prime de crédit et de liquidité selon le ténor.}
\end{solution}

\begin{solution}{iq:rc:multi-curve-and-collateral-discounting:4}
Le poseur choisira à chaque date le collatéral qui lui est le plus avantageux, le
celui qui rapporte le taux le plus élevé une fois converti dans la devise de la transaction~; les transactions
sont donc actualisées au maximum point par point des taux de collatéral~: les créances
valent moins et les dettes sont moins négatives. Avec des taux volatils et une base, l'option
possède aussi une valeur temps~; on l'évalue souvent par simulation.
\iqlookfor{la règle du maximum et le signe de l'effet.}
\end{solution}

\begin{solution}{iq:rc:multi-curve-and-collateral-discounting:5}
Les valeurs varient de la différence des valeurs actuelles des flux de trésorerie sous
les deux \omterm{def:rc:multi-curve-and-collateral-discounting:curves}{courbes d'actualisation} (le receveur de taux fixe dans la monnaie gagne lorsque les taux baissent)~;
le risque d'actualisation passe d'un indice à l'autre. Compensation équitable~: en espèces,
égale à la variation de valeur, calculée à taux forward constants, compensée par compte~;
et, lorsque les deux taux ne sont pas liés, des swaps de base qui rétablissent l'ancien profil
de risque, avec une adjudication pour les membres qui n'en veulent pas.
\iqlookfor{valeur et risque compensés séparément.}
\end{solution}

\begin{solution}{iq:rc:multi-curve-and-collateral-discounting:6}
Les swaps au pair sur l'indice propre de la \omterm{def:rc:multi-curve-and-collateral-discounting:curves}{courbe de projection} se réévaluent encore (les deux courbes sont
calibrées sur eux), de sorte que l'erreur se dissimule sur les transactions hors marché et anciennes,
où le taux d'actualisation est trop élevé de la base~: une erreur d'environ valeur
fois duration fois la base, soit 15 points de base ici, 60\,300~EUR sur un swap de dix ans
à 80 points de base dans la monnaie sur 100 millions d'euros. Un test sur une transaction hors marché
la révèle.
\iqlookfor{les transactions au pair masquent le bogue~; les transactions hors marché le révèlent.}
\end{solution}
